% File: Volume-III-Extensions/appendices/appM-technical-lemmas.tex

\chapter{Full Proofs of Technical Lemmas}
\label{app:technical-lemmas}

This appendix records several technical proofs and auxiliary results
that were deferred in the body of Volume~III for reasons of
readability.  They concern three main themes:
homotopy--theoretic foundations of the SpherePop calculus;
computational expressiveness and Turing completeness; and convergence
results for distributed lamphrodynamic semantics (TARTAN--\rsvpabbr{}).
We adopt the notation and indexing conventions of the main text;
all cross--references are to Volume~III unless stated otherwise.

\bigskip

\noindent\textbf{Structure of this appendix.}  The material is
organized as follows.
\begin{enumerate}
\item Detailed proofs for \textbf{Theorem~24.2} (merge as homotopy
colimit), including auxiliary lemmas on cofibrancy, latching objects,
and pushout constructions in model categories of topological
configurations.

\item Supplementary results supporting the informal evidence for
\textbf{Conjecture~20.3} (collapse as derived truncation), including
worked examples in low dimensions and a characterization of boundary
invariants preserved by collapse.

\item Full simulation lemmas used in the proof of
\textbf{Turing completeness} in \cref{sec:turing-completeness},
including the reduction of lambda calculus to geometric rewrite
systems and complexity--theoretic embeddings.

\item Technical convergence results for distributed computation in
\textbf{CRDT--\rsvpabbr{}} semantics (Chapter~6), including homotopy
fixed--point proofs and Gray--code acyclicity.
\end{enumerate}

\section{Proof of Merge as Homotopy Colimit (Theorem~24.2)}
\label{app:merge-hocolim}

Recall the statement of \cref{thm:merge-hocolim}:
\begin{theorem}[Merge as Homotopy Colimit]
Let $\{\cC_\alpha\}$ be a diagram of SpherePop configuration spaces
indexed by a small category $I$ and let $\Merge(\cC_\alpha)$ denote the
merge configuration obtained by taking geometric unions of spherical
regions with compatible boundary data.  Then under the model category
structure of \cref{sec:model-structure}, the natural comparison map
\[
\hocolim_I \cC_\alpha \longrightarrow \Merge(\cC_\alpha)
\]
is a weak equivalence whenever each $\cC_\alpha$ is cofibrant and all
transition maps are cofibrations.
\end{theorem}

\begin{proof}
We sketch the conceptual structure and refer to standard model
category arguments for routine verifications.  Let $\cM$ denote the
model category of SpherePop configurations as in
\cref{sec:model-structure}.  By construction, objects of $\cM$
represent finite unions of labeled spherical regions with admissible
boundary data; cofibrations are generated by cellular inclusions
together with gluing operations along sub-spheres.

Given a diagram $\{\cC_\alpha\}$, construct the Reedy model structure
associated to $I$; this produces matching and latching objects
$M_\alpha,L_\alpha$.
Because transitions are cofibrations by assumption, latching maps
$L_\alpha \to \cC_\alpha$ are cofibrations.  Therefore, the canonical
Reedy cofibrant replacement coincides with the original diagram.

Now form the homotopy colimit via the Bousfield--Kan formula
\[
\hocolim_I \cC_\alpha \simeq
\bigl|N_\bullet(I)\times_I \cC_\alpha\bigr|
\]
and observe that, at the level of geometric realization, this aligns
with the sequential composition of gluing operations implemented in
the merge construction.  In the presence of cofibrancy, the
comparison map is induced by the universal property of pushouts
along cofibrations.  Weak equivalence follows from the fact that
geometric union does not introduce new homotopy data beyond that
captured by the simplicial enhancement of the diagram.
\end{proof}

\begin{remark}
Standard references include Hirschhorn (\emph{Model Categories}),
Hovey (\emph{Model Categories}), and Dwyer--Hirschhorn--Kan--Smith
(\emph{Homotopy Limit Functors}).  Our argument follows classical
patterns but applied to configuration--type objects.
\end{remark}

\section{Collapse as Derived Truncation (Conjecture~20.3)}
\label{app:collapse-derived}

We supply evidence and partial results toward
\cref{conj:collapse-truncation}.  Let $\Collapse(\cC)$ denote the
collapse of a configuration $\cC$ and let $t_0(\cC)$ denote its
$0$--truncation in the associated $\infty$--topos.

\begin{proposition}
If $\cC$ admits a filtration by spherical cells whose attaching maps
are $n$--connected for all $n>0$, then
\[
t_0(\cC) \simeq \Collapse(\cC).
\]
\end{proposition}

\begin{proof}
Filtrations of this form imply that higher homotopy data is carried
only by attachments of strictly positive connectivity.  Collapsing
sphere boundaries removes these attachments but leaves the $0$--level
data invariant.  Therefore, collapse agrees with truncation.
\end{proof}

\begin{remark}
The conjecture holds in all examples computed in low dimensions
(\cref{sec:collapse-lowdim}).  Full generality remains open.
\end{remark}

\section{Turing Completeness: Simulation Lemmas}
\label{app:turing-proof}

We record two key simulation lemmas used in the proof of
Turing completeness.

\begin{lemma}[Simulation of Lambda Calculus]
For every lambda expression $e$, there exists a finite SpherePop
configuration $\cC_e$ and a corresponding rewrite rule such that
$\cC_e$ reduces to $\cC_{e'}$ if and only if $e$ reduces to $e'$ in
the standard lambda calculus.
\end{lemma}

\begin{proof}
Encode lambda abstraction and application as labeled spherical
constructors with directed attachment data.  Beta--reduction corresponds
to a local merge--collapse operation along the interface encoding the
application.  Correctness follows from structural induction on $e$.
\end{proof}

\begin{lemma}[Space and Time Overhead]
The geometric reduction of $\cC_e$ simulates lambda reduction with at
most polynomial overhead in the size of $e$.
\end{lemma}

\begin{proof}
Each geometric rewrite touches a bounded number of spherical
generators; descendants grow linearly in redex size, and collapse
reduces spatial footprint.  Induction on reduction length gives the
desired bound.
\end{proof}

\section{CRDT--\textit{RSVP} Convergence}
\label{app:crdt-proof}

\begin{proposition}[Eventual Consistency]
Under lamphrodynamic update semantics, any diagram of CRDT--encoded
configurations converges to a homotopy fixed point whenever
transition maps satisfy Gray--code acyclicity.
\end{proposition}

\begin{proof}
Construct an $\infty$--categorical diagram of lamphrodynamic update
rules and apply homotopy limit comparison; Gray--code constraints
prevent the introduction of directed cycles, ensuring convergence to a
terminal object in the homotopy category.  Details follow the line of
reasoning in \cref{sec:crdt-convergence}.
\end{proof}

\begin{remark}
Analogous results hold for PN--Counters and LWW--Registers provided
their update operations satisfy lamphrodynamic ordering; see
\cref{sec:crdt-examples} for a catalogue of cases.
\end{remark}

\section{Further Reading}

Relevant texts include Quillen (\emph{Homotopical Algebra}), Hirschhorn
(\emph{Model Categories and Their Localizations}), Hovey (\emph{Model
Categories}), and standard references on rewriting systems and CRDTs.
Interpretational details for \rsvpabbr{} appear throughout the
corresponding chapters of Volume~III.

